% GENERATED FILE. Edit canonical lessons, then run npm run generate:books.
% canonical-source-hash: fnv1a64:88b8c102f668fbdf
% canonical-lessons: GE-C439-stadtteil, GE-C439-viertel, GE-C439-vorort, GE-C439-polizeirevier, GE-C439-kaserne

\chapter{Stadtteile}
\label{ch:ge-stadtteile}
\hlchaptermodality{439}

\begin{chapteropening}
\textbf{By the end of this chapter:} \emph{I can talk about parts of a town: the district, the neighbourhood, the suburb, the police station and the barracks.}

Complete the last lesson of chapter 439: I can talk about parts of a town: the district, the neighbourhood, the suburb, the police station and the barracks.
\end{chapteropening}

\section[der Stadtteil]{der Stadtteil — the district, the part of town}

\label{lesson:GE-C439-stadtteil}



\begin{quote}
Before the new one: say the German for the town wall, then the German for the town gate.
\end{quote}

\subsection*{You'll want to know: der Stadtteil}
\textbf{der Stadtteil} — \textquotedblleft{}the district, the part of town\textquotedblright{}.

\begin{quote}
In unserem Stadtteil gibt es zwei Schulen.
\end{quote}

\begin{quote}
\emph{There are two schools in our part of town.}
\end{quote}

Plural: \emph{die Stadtteile}.

\begin{cousinweb}
\emph{Stadtteil} is built from \emph{Stadt} 'town' and \emph{Teil} 'part', inherited from Proto-Germanic \emph{dailą} (EtymDB).
\end{cousinweb}

\begin{grammarlens}[title={What you've built}]
One more word for telling, explaining and describing at B1.
\end{grammarlens}

\subsection*{Your turn}
\begin{itemize}
\raggedright
  \item \emph{Say it:} \emph{der Stadtteil}
  \item \emph{Say it:} \emph{der Stadtteil}, once more
  \item \emph{Say it:} a sentence of your own with \emph{der Stadtteil}
  \item \emph{Recall:} say the German for the town wall, then the German for the town gate, then say \emph{der Stadtteil} again
\end{itemize}

\begin{tcolorbox}[breakable,colback=teal!4,colframe=teal!35!black,title={Before you move on}]
What does der Stadtteil mean? (\textquotedblleft{}The district, the part of town\textquotedblright{}.) Say it once more.
\end{tcolorbox}
\section[das Viertel]{das Viertel — the neighbourhood, the quarter (of a town); a quarter}

\label{lesson:GE-C439-viertel}



\begin{quote}
Before the new one: say the German for the town gate, then the German for the district.
\end{quote}

\subsection*{You'll want to know: das Viertel}
\textbf{das Viertel} — \textquotedblleft{}the neighbourhood, the quarter (of a town); a quarter\textquotedblright{}.

\begin{quote}
Sie wohnt in einem ruhigen Viertel am Rand der Stadt.
\end{quote}

\begin{quote}
\emph{She lives in a quiet neighbourhood on the edge of town.}
\end{quote}

Same form in the plural. Also a quarter: \emph{ein Viertel nach drei}.

\begin{cousinweb}
\emph{Viertel} is built on \emph{vier} 'four' and first meant a fourth part, then a part of a town. Both are inherited Germanic words.
\end{cousinweb}

\begin{grammarlens}[title={What you've built}]
One more word for telling, explaining and describing at B1.
\end{grammarlens}

\subsection*{Your turn}
\begin{itemize}
\raggedright
  \item \emph{Say it:} \emph{das Viertel}
  \item \emph{Say it:} \emph{das Viertel}, once more
  \item \emph{Say it:} a sentence of your own with \emph{das Viertel}
  \item \emph{Recall:} say the German for the town gate, then the German for the district, then say \emph{das Viertel} again
\end{itemize}

\begin{tcolorbox}[breakable,colback=teal!4,colframe=teal!35!black,title={Before you move on}]
What does das Viertel mean? (\textquotedblleft{}The neighbourhood, the quarter (of a town); a quarter\textquotedblright{}.) Say it once more.
\end{tcolorbox}
\section[der Vorort]{der Vorort — the suburb}

\label{lesson:GE-C439-vorort}



\begin{quote}
Before the new one: say the German for the district, then the German for the neighbourhood.
\end{quote}

\subsection*{You'll want to know: der Vorort}
\textbf{der Vorort} — \textquotedblleft{}the suburb\textquotedblright{}.

\begin{quote}
Viele Familien ziehen in einen Vorort mit Garten.
\end{quote}

\begin{quote}
\emph{Many families move to a suburb with a garden.}
\end{quote}

Plural: \emph{die Vororte}.

\begin{cousinweb}
\emph{Vorort} is built from \emph{vor} 'before, in front' and \emph{Ort} 'place', inherited from Old High German \emph{ort} 'point, edge' (EtymDB).
\end{cousinweb}

\begin{grammarlens}[title={What you've built}]
One more word for telling, explaining and describing at B1.
\end{grammarlens}

\subsection*{Your turn}
\begin{itemize}
\raggedright
  \item \emph{Say it:} \emph{der Vorort}
  \item \emph{Say it:} \emph{der Vorort}, once more
  \item \emph{Say it:} a sentence of your own with \emph{der Vorort}
  \item \emph{Recall:} say the German for the district, then the German for the neighbourhood, then say \emph{der Vorort} again
\end{itemize}

\begin{tcolorbox}[breakable,colback=teal!4,colframe=teal!35!black,title={Before you move on}]
What does der Vorort mean? (\textquotedblleft{}The suburb\textquotedblright{}.) Say it once more.
\end{tcolorbox}
\section[das Polizeirevier]{das Polizeirevier — the police station}

\label{lesson:GE-C439-polizeirevier}



\begin{quote}
Before the new one: say the German for the neighbourhood, then the German for the suburb.
\end{quote}

\subsection*{You'll want to know: das Polizeirevier}
\textbf{das Polizeirevier} — \textquotedblleft{}the police station\textquotedblright{}.

\begin{quote}
Nach dem Diebstahl geht er zum Polizeirevier.
\end{quote}

\begin{quote}
\emph{After the theft he goes to the police station.}
\end{quote}

Plural: \emph{die Polizeireviere}. Also \emph{die Polizeiwache} or \emph{die Wache}.

\begin{cousinweb}
\emph{Polizeirevier} is built from \emph{Polizei} 'police', from Latin \emph{polītīa} 'the state, government' (EtymDB), and \emph{Revier} 'district', which goes back to Middle High German \emph{riviere} 'hunting ground' (EtymDB), usually traced to French.
\end{cousinweb}

\begin{grammarlens}[title={What you've built}]
One more word for telling, explaining and describing at B1.
\end{grammarlens}

\subsection*{Your turn}
\begin{itemize}
\raggedright
  \item \emph{Say it:} \emph{das Polizeirevier}
  \item \emph{Say it:} \emph{das Polizeirevier}, once more
  \item \emph{Say it:} a sentence of your own with \emph{das Polizeirevier}
  \item \emph{Recall:} say the German for the neighbourhood, then the German for the suburb, then say \emph{das Polizeirevier} again
\end{itemize}

\begin{tcolorbox}[breakable,colback=teal!4,colframe=teal!35!black,title={Before you move on}]
What does das Polizeirevier mean? (\textquotedblleft{}The police station\textquotedblright{}.) Say it once more.
\end{tcolorbox}
\section[die Kaserne]{die Kaserne — the barracks}

\label{lesson:GE-C439-kaserne}



\begin{quote}
Before the new one: say the German for the suburb, then the German for the police station.
\end{quote}

\subsection*{You'll want to know: die Kaserne}
\textbf{die Kaserne} — \textquotedblleft{}the barracks\textquotedblright{}.

\begin{quote}
Die alte Kaserne ist heute ein Museum.
\end{quote}

\begin{quote}
\emph{The old barracks is a museum today.}
\end{quote}

Plural: \emph{die Kasernen}.

\begin{cousinweb}
\emph{Kaserne} was borrowed from French \emph{caserne} (EtymDB), from Provençal \emph{cazerna}, a small room for four soldiers, which goes back to Latin \emph{quaternī} 'four at a time' (Lewis \& Short).
\end{cousinweb}

\begin{grammarlens}[title={What you've built}]
One more word for telling, explaining and describing at B1.
\end{grammarlens}

\subsection*{Your turn}
\begin{itemize}
\raggedright
  \item \emph{Say it:} \emph{die Kaserne}
  \item \emph{Say it:} \emph{die Kaserne}, once more
  \item \emph{Say it:} a sentence of your own with \emph{die Kaserne}
  \item \emph{Recall:} say the German for the suburb, then the German for the police station, then say \emph{die Kaserne} again
\end{itemize}

\begin{tcolorbox}[breakable,colback=teal!4,colframe=teal!35!black,title={Before you move on}]
What does die Kaserne mean? (\textquotedblleft{}The barracks\textquotedblright{}.) Say it once more.
\end{tcolorbox}
